% Standalone problem document, generated from the adjacent README.md.
% Compile with XeLaTeX. Mathematical content is maintained in Markdown.
\documentclass[11pt,a4paper]{article}
\usepackage[fontsize=10.5pt]{scrextend}
\usepackage[margin=22mm,headheight=15pt,headsep=9mm,footskip=11mm]{geometry}
\usepackage{fontspec}
\setmainfont{texgyrepagella}[Extension=.otf,UprightFont=*-regular,BoldFont=*-bold,ItalicFont=*-italic,BoldItalicFont=*-bolditalic]
\setsansfont{texgyreheros}[Extension=.otf,UprightFont=*-regular,BoldFont=*-bold,ItalicFont=*-italic,BoldItalicFont=*-bolditalic]
\setmonofont{texgyrecursor}[Extension=.otf,UprightFont=*-regular,BoldFont=*-bold,ItalicFont=*-italic,BoldItalicFont=*-bolditalic,Scale=MatchLowercase]
\usepackage{amsmath,amssymb}
\usepackage{unicode-math}
\setmathfont{texgyrepagella-math.otf}
\usepackage{xcolor,microtype,enumitem,needspace,fancyhdr}
\usepackage{bookmark}
\definecolor{ink}{HTML}{17344B}
\definecolor{accent}{HTML}{197D80}
\definecolor{muted}{HTML}{596773}
\hypersetup{colorlinks=true,linkcolor=accent,urlcolor=accent,pdftitle={FR-09 - Complex
equiangular tight frames with twice the
dimension},pdfauthor={Open Problems in Numerical Linear Algebra}}
\urlstyle{same}
\setlength{\parindent}{0pt}
\setlength{\parskip}{5pt plus 1pt minus 1pt}
\linespread{1.04}
\setlength{\emergencystretch}{3em}
\widowpenalty=10000
\clubpenalty=10000
\displaywidowpenalty=10000
\allowdisplaybreaks[1]
\setlist{nosep,leftmargin=1.5em}
\providecommand{\tightlist}{\setlength{\itemsep}{0pt}\setlength{\parskip}{0pt}}
\providecommand{\pandocbounded}[1]{#1}
\pagestyle{fancy}
\fancyhf{}
\fancyhead[L]{\sffamily\footnotesize\color{muted}OPEN PROBLEMS IN NUMERICAL LINEAR ALGEBRA}
\fancyhead[R]{\sffamily\bfseries\footnotesize\color{accent}FR-09}
\fancyfoot[L]{\parbox[b]{0.9\textwidth}{\sffamily\fontsize{8}{10}\selectfont\color{muted}Editorial ratings; a bounded search cannot certify that a problem remains unsolved.\\Literature check: 2026-09-10}}
\fancyfoot[R]{\sffamily\footnotesize\color{muted}\thepage}
\renewcommand{\headrulewidth}{0.3pt}
\renewcommand{\headrule}{\hbox to\headwidth{\color{accent}\leaders\hrule height \headrulewidth\hfill}}
\makeatletter
\renewcommand{\section}{\Needspace{5\baselineskip}\@startsection{section}{1}{0pt}{12pt plus 2pt minus 2pt}{4pt}{\sffamily\bfseries\large\color{ink}}}
\renewcommand{\subsection}{\Needspace{4\baselineskip}\@startsection{subsection}{2}{0pt}{9pt plus 1pt minus 1pt}{3pt}{\sffamily\bfseries\normalsize\color{ink}}}
\makeatother
\setcounter{secnumdepth}{0}
\begin{document}
{\sffamily\small\color{muted}Frames and matrix designs\par}
\vspace{3pt}
{\raggedright\hyphenpenalty=10000\exhyphenpenalty=10000\sffamily\bfseries\fontsize{21}{24}\selectfont\color{ink}Complex
equiangular tight frames with twice the dimension\par}
\vspace{7pt}
{\sffamily\small\textbf{Difficulty:} extreme\quad\textbf{Importance:} interesting
to the community\par}
{\sffamily\small\bfseries\color{accent}Status: Partially resolved\par}
\vspace{4pt}
{\color{accent}\hrule height 0.8pt}
\vspace{6pt}
\textbf{Rating rationale:} All-dimension redundancy-two ETF existence is
a difficult exact-design barrier; its direct impact is in frame theory,
sensing and reconstruction.

\textbf{Conjecture.} For every integer \(d\ge2\), there exist \(2d\)
vectors \(v_1,\ldots,v_{2d}\in\mathbb C^d\) satisfying \[
\|v_i\|_2=1,\qquad
\sum_{i=1}^{2d}v_iv_i^*=2I_d,\qquad
|v_i^*v_j|^2=\frac1{2d-1}\quad(i\ne j).
\] Thus the columns of \(V=[v_1\ \cdots\ v_{2d}]\) form a unit-norm
equiangular tight frame. The target requires every dimension, rather
than only an infinite family.

These are optimal line configurations at redundancy two: their coherence
attains the Welch bound. They provide well-balanced matrix designs for
sensing and reconstruction. This is distinct from the \(d^2\)-vector SIC
problem in FR-07.

\subsection{References}\label{references}

\begin{enumerate}
\def\labelenumi{\arabic{enumi}.}
\item
  A. Glazyrin, \emph{New constructions of optimal arrangements of \(2d\)
  lines in \(\mathbb C^d\)}, August 2026, Conjecture 1 (attributed to
  Fallon and Iverson), and Sections 2--5 for constructions.
  \href{https://arxiv.org/abs/2608.16116}{Paper}.
\item
  K. Fallon and J. W. Iverson, \emph{On the optimal arrangement of
  \(2d\) lines in \(\mathbb C^d\)}, Information and Inference 14(2)
  (2025), iaaf008.
  \href{https://doi.org/10.1093/imaiai/iaaf008}{Published paper}.
\end{enumerate}

\subsection{Status check --- 2026-09-10}\label{status-check-2026-09-10}

Glazyrin explicitly retains the all-dimension conjecture while extending
the known constructions. Searches for the Fallon--Iverson conjecture and
later redundancy-two ETF results found no full resolution. Finite
dimension verification and the new tensor/power constructions do not
cover the universal statement.

\textbf{Audit update (2026-09-10):} Rechecked Glazyrin's Conjecture 1
and construction discussion, and searched for later Fallon--Iverson
results. Added the original published conjecture source. Explicit
families provide partial cases, not every dimension. This is a bounded
literature check, not a proof that no solution exists.
\end{document}
